\begin{example}[
System prompt for generating bugs with an LM
]
\small
\texttt{$<$uploaded\_files$>$} \\
\texttt{\{\{working\_dir\}\}} \\
\texttt{$<$/uploaded\_files$>$} \\
I've uploaded a python code repository in the directory \texttt{\{\{working\_dir\}\}}. \\

Can you please install this repository?
Your goal should be to configure the repository's development environment such that existing tests pass.
You are currently in the root directory of the repository, and nothing has been installed yet.
You in an Ubuntu 22.04 environment. \\

The repository is predominantly written in Python. Here are several tips for installing it: \\

1. A good place to start is to look for a \texttt{CONTRIBUTING.[md$\vert$rst]} file, which will often contain instructions on how to install the repository and any dependencies it may have. Occasionally, the \texttt{README.md} file may also contain installation instructions. \\

2. Usually, a repository may have \texttt{setup.py} or \texttt{pyproject.toml} files which can be used to install the package. \texttt{pip install -e .} is commonly used, although many packages will also require an additional specifier that installs development packages as well (e.g. \texttt{pip install -e .[dev]}). \\

3. To check whether the repository was installed successfully, run tests and see if they pass. You can usually find tests in a \texttt{tests/} or \texttt{test/} directory. You can run tests using \texttt{pytest} or \texttt{unittest}, depending on the framework used by the repository. \\

4. Sometimes, you will need to install additional packages, often listed in a \texttt{requirements.txt} or \texttt{environment.yml} file. Also, be mindful of Ubuntu system dependencies that may need to be installed via \texttt{apt-get} (e.g. \texttt{sudo apt-get install $<$package$>$}). \\

Once you are finished with installing the repository, run the \texttt{submit} command to submit your changes for review.
\end{example}
\captionof{figure}{
Prompt provided to SWE-agent + an LM asking it to install a repository.
}
\label{fig:prompt_install_repo}